\documentclass[11pt]{article}
\usepackage[a4paper,margin=26mm]{geometry}
\usepackage[T1]{fontenc}
\usepackage{lmodern}
\usepackage{amsmath,amssymb}
\emergencystretch=2em
\title{A signed triangle remainder in a dually flat line\\Disproof of Conjecture 00000005781}
\author{ziangni-sys}
\date{4 October 2026}
\begin{document}
\maketitle
\section*{Claim and counterexample}
The original English conjecture says that the canonical divergence's
triangle deficit is a Pythagorean remainder of dual points, expressed as a
second-order interaction of potential differences, and is nonnegative.
The Chinese statement likewise asserts nonnegativity of this interaction
remainder. Neither version imposes orthogonality, projection, or an
ordering condition on the points.

Take the Euclidean line with potential $\psi(x)=x^2/2$.
Its Hessian is $\psi''=1$, so the Hessian metric is the usual positive
Euclidean metric. The primal affine coordinate is $x$ and its dual affine
coordinate is $\eta=\psi'(x)=x$. Both affine connections have zero
coefficients, so they are flat and dual for this constant metric.
Thus this is an actual dually flat structure, not an arbitrary assignment
of a divergence function.

Its Legendre-dual potential is $\varphi(\eta)=\eta^2/2$.
Indeed, completing the square gives
\[
\eta x-\psi(x)=\frac{\eta^2}{2}-\frac{(x-\eta)^2}{2}
\leq\frac{\eta^2}{2},
\]
with equality at $x=\eta$. This proves the defining supremum formula
for the Legendre transform. The canonical divergence is therefore
\[
D(x,y)=\psi(x)+\varphi(\eta(y))-x\eta(y)
=\frac{(x-y)^2}{2}.
\]
Equivalently, it is the Bregman formula
$\psi(x)-\psi(y)-\psi'(y)(x-y)$.
It is everywhere nonnegative and vanishes precisely at equal points.

\newpage
\section*{The signed three-point identity}
Define the additive triangle deficit by
\[
\Delta(x,y,z)=D(x,y)+D(y,z)-D(x,z).
\]
Expanding the squares gives the exact interaction identity
\[
\Delta(x,y,z)=(y-x)(y-z).
\]
This is the Euclidean instance of the primal/dual three-point interaction.
It has no fixed sign. For the three distinct points $(0,1,2)$,
\[
D(0,1)=\tfrac12,\quad D(1,2)=\tfrac12,\quad D(0,2)=2,
\qquad \Delta(0,1,2)=-1<0.
\]
Hence the asserted universal nonnegativity is false.
If the triangle remainder is instead defined with the opposite sign,
use the three distinct points $(0,2,1)$: their deficit is
$D(0,2)+D(2,1)-D(0,1)=2$, so the opposite remainder is $-2<0$.
The disproof thus does not depend on a choice of sign convention.
Reversing arguments in the definition of canonical divergence has no
effect here, since the quadratic example is symmetric.

Orthogonality makes this interaction zero and yields the familiar
Pythagorean equality; suitable projection conditions can make one chosen
sign nonnegative. Such hypotheses are absent from the conjecture.
The result disproves its nonnegativity clause, not the validity of a
three-point identity under its usual definitions.

\section*{Formal model and verification}
The Lean project defines the real potential and proves its actual
derivative and Hessian using Mathlib's derivative operator.
An explicit positive Jensen gap proves strict convexity.
The Legendre bound and its attaining point are proved by completing
the square. The canonical and Bregman divergences are both defined from
the potential and its derivative, and are proved equal.
The squared-distance formula, nonnegativity of the divergence itself,
and signed three-point identity follow from these definitions.
Finally, \texttt{conjecture\_5781\_false} and
\texttt{opposite\_sign\_also\_false} refute nonnegativity for both conventions.
The report's identification of the constant Hessian metric and its two
flat connections is direct coordinate geometry; no general
information-geometric theorem is assumed or claimed as formalized.
\end{document}
